\documentclass[11pt,a4paper]{article}
\usepackage[utf8]{inputenc}
\usepackage[T1]{fontenc}
\usepackage[english]{babel}
\usepackage[margin=2.5cm]{geometry}
\usepackage[dvipsnames]{xcolor}
\usepackage{amsmath,amssymb,amsfonts}
\usepackage{bm}
\usepackage{graphicx}
\usepackage{subcaption}
\usepackage{booktabs}
\usepackage{float}
\usepackage{siunitx}
\usepackage[colorlinks,linkcolor=blue,citecolor=blue,urlcolor=blue]{hyperref}
%
% Minimal aliases, taken from NEGF-papers/common/headers/{alias,math}.tex
%
% Colors (notes to self, remove before submission)
\newcommand{\note}[1]{\textcolor{blue}{[#1]}}
\newcommand{\todo}[1]{\textcolor{violet}{[#1]}}
%
% Real space
\def\ReS  {\mathbb{R}}
\newcommand{\di}{\mathrm{d}}
%
% Equations
\newcommand{\eq}[1]{\begin{align}#1\end{align}}
\newcommand{\eqn}[1]{\begin{align*}#1\end{align*}}
\newcommand{\eql}[2]{\begin{align}\label{#1}#2\end{align}}
\newcommand{\seqal}[2]{\begin{subequations}\label{#1}\begin{align}#2\end{align}\end{subequations}}
%
% Functions
\newcommand{\average}[1]{\left\langle #1 \right\rangle}
\newcommand{\pr}[1]{\left( #1 \right)}
\newcommand{\ps}[1]{\left[ #1 \right]}
\newcommand{\lab}[1]{\label{#1}}
\newcommand{\bs}[1]{\boldsymbol{#1}}
\newcommand{\e}[1]{Eq.~\eqref{#1}}
\newcommand{\es}[2]{Eqs.~\eqref{#1}--\eqref{#2}}
\newcommand{\tab}[1]{Tab.~\ref{#1}}
\newcommand{\fig}[1]{Fig.~\ref{#1}}
\renewcommand{\sec}[1]{Section~\ref{#1}}
%
% ML notation
\newcommand{\pd}[2]{\frac{\partial #1}{\partial #2}}
\def\bth        {\bs{\theta}}
\def\xI         {x^{(i)}}
\def\yI         {y^{(i)}}
\def\hth        {h_{\bth}}


\title{Intelligent Systems -- Machine Learning (MICS2-62)\\[2pt]
\large Take-Home Assignment 1}
\author{Riccardo Reho\\ \small individual submission}
\date{\today}

\begin{document}
\maketitle

\section{A. Supervised learning: regression}
\lab{sec:A}

\subsection{A.I Linear regression}
\lab{sec:A1}

\paragraph{A.I.1 Data acquisition.}
I am not yet assigned to a group as I am an external student (Postdoctoral Researcher at UniLu). 
As an individual submission I take the data block of group $n=1$, i.e.\ Excel rows
$(n-1)\cdot 20+2=2$ to $n\cdot 20+1=21$. The resulting training set is
\eql{eq:A1_dataset}{
\mathcal{D}=\left\{\pr{\xI,\yI}\right\}_{i=1}^{m},\qquad m=20,\qquad \xI\in\ReS\ ,\quad \yI\in\ReS\ ,
}
with $\xI$ the ambient temperature $[^\circ\mathrm{C}]$ and $\yI$ the net 
hourly electrical energy output $[\mathrm{MW}]$, listed in \tab{tab:A1_data}.

\noindent\emph{Code} (\texttt{code/data.py}). The class \texttt{PowerPlantData(n)} reads the Excel file with
\texttt{pandas} and, working with 0-based indices, selects the $m$ consecutive data points starting at
$\texttt{start}=(n-1)\,m$, i.e.\ \texttt{df.iloc[start:start+m]}. Since row~1 of the Excel file is the header and
Excel counts from 1, data point \texttt{df.iloc[k]} sits in Excel row $k+2$; this conversion is used only to label
\tab{tab:A1_data}. The selected points are stored as two \texttt{numpy} arrays \texttt{x}, \texttt{y}, and
\texttt{to\_latex} writes \tab{tab:A1_data}.

\begin{table}[H]
\centering
\caption{Data set $\mathcal{D}$ of \e{eq:A1_dataset}: Excel rows 2--21.}
\lab{tab:A1_data}
\input{tables/A1_data}
\end{table}

\paragraph{A.I.2 Data transformation.}
I apply $\mathrm{Z}$-score standardization to both the input and the output,
\eql{eq:A2_zscore}{
\tilde x^{(i)}=\frac{\xI-\mu_x}{\sigma_x},\qquad
\tilde y^{(i)}=\frac{\yI-\mu_y}{\sigma_y},\qquad
\mu_v=\frac{1}{m}\sum_{i=1}^{m}v^{(i)},\quad
\sigma_v=\sqrt{\frac{1}{m}\sum_{i=1}^{m}\pr{v^{(i)}-\mu_v}^2},
}
with $v\in\{x,y\}$. The transformed data have zero mean and unit variance; they are listed in \tab{tab:A2_zscore}
together with $\mu_v$ and $\sigma_v$.

\noindent\emph{Justification.} From the data I can infer that $x\sim 20$ and $y\sim 450$, 
so the two components of the gradient
of the cost function (\sec{sec:A1}) differ by orders of 
magnitude and gradient descent would need a very small
learning rate to remain stable. After standardization both 
variables are $\mathcal{O}(1)$, the cost function is
well conditioned and a single learning rate works for all parameters. 
I prefer the $\mathrm{Z}$-score to Min-Max
normalization because it centres the data. The least-squares 
intercept then vanishes, $\tilde\theta_0=0$, which
gives an analytic check on the results. Standardizing $y$ as
well keeps the cost function $\mathcal{O}(1)$. The transformation is linear and invertible,
$\yI=\mu_y+\sigma_y\tilde y^{(i)}$.

\noindent\emph{Code} (\texttt{code/transform.py}). The class \texttt{ZScore} implements \e{eq:A2_zscore}:
\texttt{fit} computes $\mu$ and $\sigma$ (population standard deviation, \texttt{ddof=0}), \texttt{transform}
returns $(v-\mu)/\sigma$ and \texttt{inverse} maps back to physical units. One instance is fitted on $x$ and one
on $y$.

\begin{table}[H]
\centering
\caption{Data of \tab{tab:A1_data} before and after the z-score standardization of \e{eq:A2_zscore}.}
\lab{tab:A2_zscore}
\input{tables/A2_zscore}
\end{table}

\paragraph{A.I.3 Least squares.}
\input{tables/A3_theta}%
I fit the linear model $\hth(\tilde x)=\theta_0+\theta_1\tilde x$ to the standardized data of \tab{tab:A2_zscore}.
Collecting the data in the design matrix $X\in\ReS^{m\times 2}$, whose $i$-th row is $(1,\tilde x^{(i)})$, and the
outputs in $\tilde{\bs y}=(\tilde y^{(1)},\dots,\tilde y^{(m)})^T$, least squares minimizes
$\|X\bth-\tilde{\bs y}\|^2=\sum_i r_i^2$, with residuals $\bs r=X\bth-\tilde{\bs y}$, i.e.\
$r_i=\sum_k X_{ik}\theta_k-\tilde y^{(i)}$. Since $\partial r_i/\partial\theta_j=X_{ij}$, the chain rule gives
\eql{eq:A3_grad}{
\pd{}{\theta_j}\sum_i r_i^2=2\sum_i X_{ij}\,r_i=2\pr{X^T\bs r}_j
\qquad\Longrightarrow\qquad
\nabla_{\bth}\|X\bth-\tilde{\bs y}\|^2=2X^T\pr{X\bth-\tilde{\bs y}}.
}
Geometrically, $X\bth$ spans the plane generated by the two columns of $X$ in $\ReS^m$, and the minimum is reached
when $X\bth$ is the orthogonal projection of $\tilde{\bs y}$ onto it, i.e.\ when the residual is orthogonal to both
columns, $X^T\bs r=0$. Setting \e{eq:A3_grad} to zero gives the normal equations
\eql{eq:A3_normal}{
X^TX\,\bth=X^T\tilde{\bs y}
\quad\Longleftrightarrow\quad
\begin{pmatrix} m & \sum_i\tilde x^{(i)} \\ \sum_i\tilde x^{(i)} & \sum_i\pr{\tilde x^{(i)}}^2\end{pmatrix}
\begin{pmatrix}\theta_0\\ \theta_1\end{pmatrix}
=\begin{pmatrix}\sum_i\tilde y^{(i)}\\ \sum_i\tilde x^{(i)}\tilde y^{(i)}\end{pmatrix}.
}
The first row gives $\theta_0=\average{\tilde y}-\theta_1\average{\tilde x}$: the least-squares line always passes
through the centroid of the data. After standardization $\sum_i\tilde x^{(i)}=\sum_i\tilde y^{(i)}=0$ and
$\sum_i(\tilde x^{(i)})^2=m$, so the centroid is the origin and \e{eq:A3_normal} reduces to
\eql{eq:A3_solution}{
\theta_0=0,\qquad \theta_1=\frac{1}{m}\sum_{i=1}^m\tilde x^{(i)}\tilde y^{(i)},
}
i.e.\ $\theta_1$ is the Pearson correlation coefficient of $x$ and $y$. Numerically I obtain
\eqn{
\theta_0=\num{\LSthO}\simeq 0,\qquad \theta_1=\LSthI,
}
with $\theta_0$ zero to machine precision, as expected. Inverting \e{eq:A2_zscore} gives the model in physical units,
\eql{eq:A3_phys}{
y&=\theta_0^{\mathrm{ph}}+\theta_1^{\mathrm{ph}}x, \nonumber\\
\theta_1^{\mathrm{ph}}&=\frac{\sigma_y}{\sigma_x}\theta_1=\LSphI\ \mathrm{MW}/^\circ\mathrm{C}, \nonumber\\
\theta_0^{\mathrm{ph}}&=\mu_y+\sigma_y\theta_0-\theta_1^{\mathrm{ph}}\mu_x=\LSphO\ \mathrm{MW}:
}
the net output drops by about $2.2$\,MW per degree of ambient temperature. The fit is shown in \fig{fig:A3}.

\noindent\emph{Code} (\texttt{code/least\_squares.py}). The class \texttt{LeastSquares} builds $X$ and solves
\e{eq:A3_normal} with \texttt{np.linalg.solve} (which avoids computing $(X^TX)^{-1}$ explicitly); \texttt{predict}
evaluates $\hth$. The function \texttt{to\_physical} implements \e{eq:A3_phys}.

\begin{figure}[H]
\centering
\includegraphics[width=\linewidth]{figs/A3_least_squares.pdf}
\caption{Least-squares regression line and data points, in standardized (left) and physical (right) units.}
\lab{fig:A3}
\end{figure}

\paragraph{A.I.4 Gradient descent: cost function.}
For gradient descent I keep the model $\hth(\tilde x)=\theta_0+\theta_1\tilde x$ and choose the mean squared error
\eql{eq:A4_cost}{
J(\bth)=\frac{1}{2m}\sum_{i=1}^{m}\pr{\hth(\tilde x^{(i)})-\tilde y^{(i)}}^2=\frac{1}{2m}\|X\bth-\tilde{\bs y}\|^2 .
}
\emph{Motivation.} (i) $J$ is proportional to the least-squares objective of A.I.3, so it has the same unique
minimizer and the two methods can be compared directly. (ii) $J$ is a convex quadratic function of $\bth$: it has no
local minima other than the global one, and gradient descent converges to it for a small enough learning rate.
(iii) It is differentiable everywhere, with a gradient linear in the residuals. (iv) The factor $1/m$ makes $J$
independent of the number of data points, and the factor $1/2$ cancels the 2 coming from the derivative of the square.

\noindent\emph{Partial derivatives.} With $r_i=\hth(\tilde x^{(i)})-\tilde y^{(i)}$,
\seqal{eq:A4_grad}{
\pd{J}{\theta_0}&=\frac{1}{m}\sum_{i=1}^{m}\pr{\hth(\tilde x^{(i)})-\tilde y^{(i)}},\\
\pd{J}{\theta_1}&=\frac{1}{m}\sum_{i=1}^{m}\pr{\hth(\tilde x^{(i)})-\tilde y^{(i)}}\tilde x^{(i)},
}
i.e.\ $\nabla_{\bth}J=\frac{1}{m}X^T\pr{X\bth-\tilde{\bs y}}$, the same gradient as \e{eq:A3_grad} up to a factor.

\noindent\emph{Update rules.} At each iteration $t$ both parameters are updated simultaneously with learning rate
$\alpha>0$,
\seqal{eq:A4_update}{
\theta_0^{(t+1)}&=\theta_0^{(t)}-\alpha\,\pd{J}{\theta_0}\Big|_{\bth^{(t)}},\\
\theta_1^{(t+1)}&=\theta_1^{(t)}-\alpha\,\pd{J}{\theta_1}\Big|_{\bth^{(t)}},
}
or compactly $\bth^{(t+1)}=\bth^{(t)}-\frac{\alpha}{m}X^T\pr{X\bth^{(t)}-\tilde{\bs y}}$. A fixed point of
\e{eq:A4_update} satisfies $X^T(X\bth-\tilde{\bs y})=0$, which is exactly \e{eq:A3_normal}.

\noindent\emph{Code} (\texttt{code/gradient\_descent.py}). The class \texttt{LinearGD} stores $\bth$ and $\alpha$;
\texttt{cost} implements \e{eq:A4_cost}, \texttt{gradient} implements \e{eq:A4_grad},
and \texttt{step} applies \e{eq:A4_update}, recording $\bth^{(t)}$ and $J(\bth^{(t)})$ at each
iteration.

\paragraph{A.I.5 Gradient descent: first iteration.}
\input{tables/A5_values}%
I draw the initial parameters uniformly in $[-1,1]$ with a seeded random generator (seed \GDseed, for
reproducibility), which gives
\eqn{
\theta_0^{(0)}=\GDthOinit,\qquad \theta_1^{(0)}=\GDthIinit,
}
and use the learning rate $\alpha=\GDalpha$. Since the data are standardized the curvature of $J$ is $\mathcal{O}(1)$
(see A.I.8), so any $\alpha<2$ converges; I take a small value so that the approach to the minimum is visible over
several iterations. The first iteration evaluates the gradient \e{eq:A4_grad} at $\bth^{(0)}$ and applies
\e{eq:A4_update}; the computed values are in the first row of \tab{tab:A5_iterations} and the resulting regression
line is in \fig{fig:A5_fits} (left).

\begin{table}[H]
\centering
\caption{First two iterations of gradient descent, \e{eq:A4_update}, with $\alpha=\GDalpha$: parameters before the
update, gradient \e{eq:A4_grad}, parameters after the update and cost \e{eq:A4_cost}.}
\lab{tab:A5_iterations}
\small
\input{tables/A5_iterations}
\end{table}

\paragraph{A.I.6 Gradient descent: second iteration.}
The second iteration starts from $\bth^{(1)}$; its values are in the second row of \tab{tab:A5_iterations} and the
regression line in \fig{fig:A5_fits} (centre). The cost decreases at each step, and the line rotates towards the
data.

\paragraph{A.I.7 Gradient descent: last iteration.}
I iterate \e{eq:A4_update} until the gradient norm drops below $\num{\GDtol}$, which happens after $T=\GDT$ iterations.
The final parameters are
\eqn{
\theta_0^{(T)}=\num{\GDthO},\qquad \theta_1^{(T)}=\GDthI,\qquad J(\bth^{(T)})=\GDJ,
}
corresponding through \e{eq:A3_phys} to $\theta_0^{\mathrm{ph}}=\GDphO$\,MW and
$\theta_1^{\mathrm{ph}}=\GDphI$\,MW/$^\circ$C. They agree with the least-squares solution of A.I.3 to within
$\num{\GDdiff}$. The final regression line is in \fig{fig:A5_fits} (right) and the cost for all iterations in
\fig{fig:A7_cost}.

\begin{figure}[H]
\centering
\includegraphics[width=\linewidth]{figs/A5_fits.pdf}
\caption{Data and regression line after the first (A.I.5), second (A.I.6) and last (A.I.7) iteration of gradient
descent. The dashed line is the least-squares solution of A.I.3.}
\lab{fig:A5_fits}
\end{figure}

\begin{figure}[H]
\centering
\includegraphics[width=0.9\linewidth]{figs/A7_cost.pdf}
\caption{Cost function $J(\bth^{(t)})$ for all iterations (left) and its excess over the minimum on a logarithmic
scale (right): the straight line shows the geometric convergence of \e{eq:A8_contraction}.}
\lab{fig:A7_cost}
\end{figure}

\noindent\emph{Code} (\texttt{code/gd\_linear.py}). The script fits the scalers, draws $\bth^{(0)}$, runs
\texttt{LinearGD.run} with the stopping criterion $\|\nabla_{\bth}J\|<\num{\GDtol}$, and writes
\tab{tab:A5_iterations} and Figs.~\ref{fig:A5_fits}--\ref{fig:A7_cost} from the stored history.

\paragraph{A.I.8 Discussion.}
Least squares and gradient descent minimize the same cost: $J$ of \e{eq:A4_cost} is the least-squares objective
divided by $2m$. They therefore share the same model, the same optimum and, once converged, the same regression line
(\fig{fig:A5_fits}, right). The difference lies in how the minimum is reached.

\emph{Least squares} solves the stationarity condition $\nabla_{\bth}J=0$ exactly, in one step, by solving the
$2\times2$ linear system \e{eq:A3_normal}. It has no hyperparameters and gives the minimum to machine precision
($\theta_0\sim10^{-16}$). Its cost is dominated by forming and solving $X^TX$, i.e.\ $\mathcal{O}(mp^2+p^3)$ for $p$
parameters, which becomes expensive for many features; it exists only when $J$ is quadratic (linear model and squared
error); and it fails when $X^TX$ is singular or badly conditioned.

\emph{Gradient descent} reaches the minimum iteratively. It requires choosing $\theta^{(0)}$, $\alpha$ and a stopping
criterion, and it stops at a finite tolerance ($\theta_0\sim10^{-7}$ here). Each iteration costs only
$\mathcal{O}(mp)$, no matrix has to be inverted, and the same update rule applies to any differentiable cost, which is
exactly what A.II needs (a quartic cost has no closed-form solution).

\emph{Computational cost.} With $X\in\ReS^{m\times p}$ ($p=2$ here), least squares requires: forming $X^TX$, whose
$p^2$ entries are each a sum over the $m$ data points, $\mathcal{O}(mp^2)$; forming $X^T\tilde{\bs y}$, $p$ sums of
$m$ terms, $\mathcal{O}(mp)$; and solving the $p\times p$ system by Gaussian elimination (LU), $\mathcal{O}(p^3)$
($\simeq\frac23p^3$ operations). One gradient-descent iteration requires $X\bth$ ($m$ dot products of length $p$,
$\mathcal{O}(mp)$), the residuals ($\mathcal{O}(m)$) and $X^T\bs r$ ($p$ sums over $m$, $\mathcal{O}(mp)$), so
$\mathcal{O}(mp)$ per iteration and $\mathcal{O}(Tmp)$ in total. Gradient descent is cheaper when $T\ll p$, i.e.\
for many features and few iterations; for $p=2$ both costs are negligible.

For the standardized data the convergence can be predicted exactly. Since $\sum_i\tilde x^{(i)}=0$ and
$\sum_i(\tilde x^{(i)})^2=m$, the Hessian of $J$ is the identity, $X^TX/m=I$, so $\nabla_{\bth}J=\bth-\bth_{\mathrm{LS}}$
(as visible in \tab{tab:A5_iterations}: $\partial_{\theta_0}J=\theta_0$). The update \e{eq:A4_update} then becomes
\eql{eq:A8_contraction}{
\bth^{(t+1)}-\bth_{\mathrm{LS}}=(1-\alpha)\pr{\bth^{(t)}-\bth_{\mathrm{LS}}},\qquad
J(\bth^{(t)})-J_{\min}=\tfrac12(1-\alpha)^{2t}\,\|\bth^{(0)}-\bth_{\mathrm{LS}}\|^2 .
}
The error shrinks by a factor $1-\alpha=0.9$ per iteration (the straight line in \fig{fig:A7_cost}, right), and the
predicted number of iterations to reach the tolerance,
$\ln\!\pr{\num{\GDtol}/\|\bth^{(0)}-\bth_{\mathrm{LS}}\|}/\ln(1-\alpha)$, matches the observed $T=\GDT$. The same formula
shows the role of $\alpha$: $\alpha=1$ would reach the least-squares solution in a single step, and $\alpha>2$ would
diverge. Without standardization the Hessian would have eigenvalues differing by orders of magnitude, and a single
$\alpha$ small enough for the steep direction would make convergence along the flat one very slow: this is the
practical reason for the transformation of A.I.2.


\subsection{A.II Polynomial regression}
\lab{sec:A2}
\todo{}

\section{B. Supervised learning: classification}
\lab{sec:B}
\subsection{B.I Classification models comparison}
\todo{}
\subsection{B.II Bias--variance tradeoff}
\todo{}

\section{C. Unsupervised learning}
\lab{sec:C}
\todo{}

\section{D. Model evaluation and validation}
\lab{sec:D}
\todo{}

\end{document}
